\clearpage\begingroup\SourceLeaf{381}\ArithmeticHead{50}{L}{STÄCKEL, GAUSS ALS GEOMETER.}
\renewcommand{\thefootnote}{\arabic{footnote})}\setcounter{footnote}{0}
\begin{center}
\textbf{12.}\par\smallskip
\textls[100]{Verknotungen und Verkettungen von Kurven.}
\end{center}
\noindent\hspace*{\parindent}Eine der ältesten Aufzeichnungen von \textsc{Gauss}, die uns überhaupt im Nachlass
erhalten sind, ist ein Blatt mit der Jahreszahl 1794. Es trägt die Überschrift:
\textls[100]{A collection of knots} und enthält 13 sauber gezeichnete Ansichten
von Knoten mit daneben geschriebenen englischen Namen; man darf wohl
annehmen, dass es sich um einen Auszug aus einem englischen Buche über
Knoten handelt. Dabei liegen zwei weitere Zettel mit Zeichnungen von
Knoten; der eine ist datiert 1819, der andere stammt wohl aus noch späterer
Zeit, denn \textsc{Gauss} hat darauf vermerkt: »\textsc{Riedl}, \textit{Beiträge zur Theorie des Sehnenwinkels},
Wien 1827«.
\par
Auf die Verknotungen geschlossener Kurven beziehen sich die Bemerkungen,
die aus dem Nachlass W. VIII, S. 271—285 abgedruckt sind. Im Besonderen
hat \textsc{Gauss} in einer aus dem Dezember 1844 stammenden Notiz die zahlreichen
Formen ermittelt, die geschlossene Kurven mit vier Knoten aufweisen können.
\par
Die Verkettung von zwei Kurven im Raume betrifft die schon erwähnte
Bemerkung vom 22. Januar 1833 (W. V, S. 605), in der am Schluss die bekannte
Integralformel für die Anzahl der Umschlingungen mitgeteilt wird.
»Es war damit der erste Anfang gemacht worden zu der später vor allem
durch die von W. \textsc{Dyck} benutzte \textsc{Kronecker}sche Charakteristikentheorie erfolgreichen
Anwendung der höheren Analysis auf die Geometria situs«\footnote{M. \textsc{Dehn} und P. \textsc{Heegaard}, a. a. O., S. 155. Man findet hier auch ausführliche Angaben über die anschliessenden Arbeiten. Hinzuzufügen ist, dass \textsc{Fr. Zöllner}, \textit{Naturwissenschaft und christliche Offenbarung}, Leipzig 1881, S. 100 berichtet, ein gewisser \textsc{Schürlein}, ein Schüler von \textsc{Gauss}, habe sich sehr eingehend und unter stetiger Teilnahme von \textsc{Gauss} mit diesem Gegenstande beschäftigt; leider ist es nicht möglich gewesen, Näheres hierüber zu ermitteln.}.
\par
Die Bestimmung der gegenseitigen Lage von Kurven in der Ebene ist
das Mittel, dessen sich \textsc{Gauss} in seiner Dissertation (1799) bei der Herleitung
des Fundamentalsatzes der Algebra bedient hatte\footnote{In der Fussnote zum art. 21 der Dissertation sagt \textsc{Gauss} ausdrücklich, Beweise, die sich auf die Geometria situs stützten, seien nicht weniger schlüssig als solche, bei denen man sich der Prinzipien der Geometria magnitudinis bediene.}. Noch stärker tritt dieser
Gesichtspunkt bei der neuen Darstellung vom Jahre 1849 hervor: »Ich werde
die Beweisführung in einer der Geometrie der Lage entnommenen Einkleidung
darstellen, weil jene dadurch die grösste Anschaulichkeit und Einfachheit
\clearpage\SourceLeaf{382a}\ArithmeticHead{51}{R}{GEOMETRIA SITUS.}\setcounter{footnote}{0}
\noindent gewinnt« (W. III, S. 79). Es folgt die vorher (S. 49) angeführte Bemerkung
über die nach der Stetigkeit zusammenhängenden Grössenkombinationen. Hierin
liegt jedoch keine Einschränkung, weil »zwar die räumliche Anschauung der
beste Führer in der Entdeckung neuer Sätze [der Geometria situs] und ihrer
Beweise ist, man aber in jedem einzelnen dieser Fälle sehen kann, dass die
in Betracht kommenden Schlüsse auch allein mit Hilfe abstrakter Entwicklungen
gemacht werden können«\footnote{M. \textsc{Dehn} und P. \textsc{Heegaard}, a. a. O., S. 170.}.
\par
Endlich sind noch die Untersuchungen zu nennen, die \textsc{Gauss} über die
möglichen Verteilungsarten der geozentrischen Örter eines Planeten auf dem
Zodiakus angestellt hat (W. VI, S. 106), und die hierbei erwähnten Fälle eines
kettenartigen Ineinandergreifens zweier Planetenbahnen, wie es bei den Asteroiden
mehrfach verwirklicht ist.
\clearpage\endgroup

